\documentclass[11pt,a4paper]{article}
\usepackage[margin=2.5cm]{geometry}
\usepackage{amsmath,amssymb}

\newcommand{\R}{\mathbb{R}}
\newcommand{\E}{\mathbb{E}}
\newcommand{\norm}[1]{\left\| #1 \right\|}
\newcommand{\ip}[2]{\left\langle #1, #2 \right\rangle}

\title{Homework 19.10.26 -- Exercises 20, 24, 25}
\author{Gor Melkumyan}
\date{}

\begin{document}
\maketitle

\section*{Exercise 20}

Given that $f$ has $L$-Lipschitz gradient, i.e.,
\[
  \norm{\nabla f(x) - \nabla f(y)} \le L\norm{x-y} \qquad \forall x, y \in \R^d .
\]
Show that
\[
  -\frac{L}{2}\norm{x-y}^2 \le D_f(x,y) \le \frac{L}{2}\norm{x-y}^2 \qquad \forall x, y \in \R^d .
\]

\subsection*{Proof}
It is enough to show that $|D_f(x,y)| \le \frac{L}{2}\norm{x-y}^2$.
Let $z := y + t(x-y)$, so that $x - y = \frac{z-y}{t}$. Then
\begin{align*}
  |D_f(x,y)|
  &= \bigl|f(x) - f(y) - \ip{\nabla f(y)}{x-y}\bigr|
  \\
  &= \left|\int_0^1 \ip{\nabla f\bigl(y + t(x-y)\bigr)}{x-y}\,dt - \ip{\nabla f(y)}{x-y}\right| \\
  &= \left|\int_0^1 \ip{\nabla f(z) - \nabla f(y)}{z-y}\cdot\frac{1}{t}\,dt\right|
  \le \int_0^1 \frac{1}{t}\,\bigl|\ip{\nabla f(z) - \nabla f(y)}{z-y}\bigr|\,dt .
\end{align*}
By Cauchy--Schwarz and then by the $L$-Lipschitz gradient assumption,
\begin{align*}
  |D_f(x,y)|
  &\le \int_0^1 \frac{1}{t}\norm{z-y}\norm{\nabla f(z) - \nabla f(y)}\,dt
  \le \int_0^1 \frac{L}{t}\norm{z-y}^2\,dt
  = \int_0^1 \frac{L}{t}\ip{z-y}{z-y}\,dt \\
  &= \int_0^1 L\ip{x-y}{x-y}\,t\,dt
  = \frac{L}{2}\norm{x-y}^2,
\end{align*}
where we used $z - y = t(x-y)$.

\section*{Exercise 24}

For all $a, b \in \R^d$ and $t > 0$ the following inequalities hold:
\begin{align*}
  &1)\quad \ip{a}{b} \le \frac{\norm{a}^2}{2t} + \frac{t\norm{b}^2}{2}, \\
  &2)\quad \norm{a+b}^2 \le 2\norm{a}^2 + 2\norm{b}^2, \\
  &3)\quad \frac{1}{2}\norm{a}^2 - \norm{b}^2 \le \norm{a+b}^2 .
\end{align*}

\subsection*{Proof of 1)}
Let $a' = \frac{a}{\sqrt{t}}$, $b' = b\sqrt{t}$, i.e., $a = a'\sqrt{t}$, $b = \frac{b'}{\sqrt{t}}$. Then
\[
  \ip{a}{b} = \ip{a'}{b'}, \qquad
  \frac{\norm{a}^2}{2t} = \frac{\norm{a'}^2}{2}, \qquad
  \frac{t\norm{b}^2}{2} = \frac{\norm{b'}^2}{2}.
\]
As a result, we only have to prove
\[
  2\ip{a'}{b'} \le \norm{a'}^2 + \norm{b'}^2,
\]
which is equivalent to $\norm{a'-b'}^2 \ge 0$, so it is true.

\subsection*{Proof of 2)}
Using $\norm{a-b}^2 \ge 0$ we get
\[
  \norm{a}^2 + \norm{b}^2 - 2\ip{a}{b} \ge 0 .
\]
Adding $\norm{a}^2 + \norm{b}^2$ to both sides:
\[
  2\norm{a}^2 + 2\norm{b}^2 - 2\ip{a}{b} \ge \norm{a}^2 + \norm{b}^2 .
\]
Adding $2\ip{a}{b}$ and regrouping:
\[
  2\norm{a}^2 + 2\norm{b}^2 \ge \norm{a}^2 + 2\ip{a}{b} + \norm{b}^2 = \norm{a+b}^2 .
\]

\subsection*{Proof of 3)}
3) is the same as
\[
  \norm{a}^2 \le 2\norm{a+b}^2 + 2\norm{b}^2 .
\]
Setting $a := a' + b'$, $b := -b'$ (so that $a + b = a'$) we get
\[
  \norm{a'+b'}^2 \le 2\norm{a'}^2 + 2\norm{b'}^2,
\]
which is the same as 2).

\section*{Exercise 25}

For a random vector $X \in \R^d$ (with finite second moment) and any $c \in \R^d$,
\[
  \E\bigl[\norm{X - \E[X]}^2\bigr] = \E\bigl[\norm{X - c}^2\bigr] - \norm{\E[X] - c}^2 .
\]

\subsection*{Proof}
\begin{align*}
  \E\bigl[\norm{X - \E[X]}^2\bigr]
  &= \E\bigl[\norm{(X - c) - (\E[X] - c)}^2\bigr] \\
  &= \E\bigl[\norm{X - c}^2\bigr] + \E\bigl[\norm{\E[X] - c}^2\bigr]
     - 2\,\E\bigl[\ip{X-c}{\E[X]-c}\bigr].
\end{align*}
On the other hand, $\norm{\E[X] - c}^2$ is a number, so its expectation is the same number:
\[
  \E\bigl[\norm{\E[X] - c}^2\bigr] = \norm{\E[X] - c}^2 .
\]
So we get
\[
  \E\bigl[\norm{X - \E[X]}^2\bigr]
  = \E\bigl[\norm{X - c}^2\bigr] - \norm{\E[X] - c}^2
    - 2\Bigl(\E\bigl[\ip{X-c}{\E[X]-c}\bigr] - \norm{\E[X] - c}^2\Bigr).
\]
Also
\begin{align*}
  \E\bigl[\ip{X-c}{\E[X]-c}\bigr]
  &= \E\bigl[\ip{X}{\E[X]}\bigr] - \ip{c}{\E[X]} - \ip{c}{\E[X]} + \norm{c}^2 \\
  &= \norm{\E[X]}^2 - 2\ip{c}{\E[X]} + \norm{c}^2 = \norm{\E[X] - c}^2,
\end{align*}
which means
\[
  -2\Bigl(\E\bigl[\ip{X-c}{\E[X]-c}\bigr] - \norm{\E[X] - c}^2\Bigr) = 0,
\]
which proves the identity.

\end{document}
